\documentclass{article}
\usepackage[margin=1in]{geometry}
\usepackage[skip=1em, indent=12pt]{parskip}
\usepackage{amsmath}
\usepackage{amssymb}
\usepackage{graphicx}
\usepackage{microtype}
\usepackage{textcomp, gensymb}
\begin{document}
\setcounter{section}{5}
\section{Chapter 7: Robustness} 
\setcounter{subsection}{1}
\subsection{Variability}
Despite our best estimates, the performance of CMOS circuits can vary as a function of process, environment (i.e., temperature), and aging. 

Voltage variation can occur as a result of higher switching frequencies, and thus current, in certain areas of a circuit. This surge in current demand can cause voltage droop. Thus, \(V_{DD}\) is usually designed with \(\pm 10\%\). Ambient temperature on chip can increase as a result of chip power consumption. Ambient temperature variation ratings differ depending on application. Threshold voltage depends on the placement of the \textit{placement} of dopants in the the channel, though this process has some randomness built into it. 

Spatial distribution variation can affect transistor operation in a measurable way as well (look into this one more, I zoned out). There is also parameter variation in process, where threshold voltage, gate oxide thickness, and the effective channel length will slow down or speed up due to these parameters.

High \(V_{DD}\) will speed up a transistor's operation, and lower temperatures will have the same effect on propagation delay. Using each of these parameters and their extremes, we can use the corners of a "square" generated by the possible variation to target the general operating regions of the circuit. "Corner checking" uses a matrix to narrow down quantitative performance limits.

Critical simulation corners include cycle time, power, and subthreshold leakage.

Review of TSMC cell library, using datasheet to derive design choices. Review this stuff again.

Monte Carlo simulation is often used to determine the average impact of some parameters by running random parameters through multiple simulations.

\subsection{Reliability}
Bathtub curve; failure rate is high at first (infant mortality) and then increases at the end of its useful life. MTBF can be used to quantify this on average, and FIT.

Hot carriers. Electric fields may impart high energies to charge carriers, which become trapped in the oxide. This causes a shift in \(V_t\) over time. Choose \(V_{DD}\) to achieve a reasonable product lifetime.

NBTI (negative bias temperature instability) causes dangling bonds that traps charge carriers and causes a shift in \(V_t\). 

Accelerated lifetime testing uses high temperature or other conditions to create a more stressful environment for the chips, which causes them to wear out faster. "Electron wind" causes the movement of metal atoms along wires, and this can lead to open circuits in extreme cases. This happens much more often in DC, and electromigration is more of a risk than in AC circuits. This can cause erosion of vias, for example, as well.

Self-heating is when a wire heats itself up from current and then resistance increases. Overvoltage can be caused by ESD, static electricity discharges can be in the kilovolt range. 

Random bit flips can occur due to alpha, beta radiation, solar radiation.

Latchup occurs when a pnp transistor is generated by current flowing through the substrate (a parasitic transistor). This shorts \(V_{DD}\) and \(GND\) and destroys the device as a result. External voltages on a PCB can be unpredictable and/or much larger than the CMOS circuit, so I/O pins must be carefully managed as they are the connection between CMOS chips and a larger PCB environment. Guard rings are used to handle this. 

\subsection{Scaling}
Scale decreases exponentially every year in VLSI (Dennard's scaling). MOSFET scaling has knock-on effects on \(\beta\), gate capacitance, clock frequency, switching energy, power, area per gate, and switching current density. Real scaling has slowed from Dennard scaling, as oxide thicknesses reach the atomic scale, and \(V_{DD}\) scaling has slowed due to its impact on threshold effects.

Interconnect scaling also has an enormous impact on CMOS functionality.

\end{document}
